\honest{Policy Debates Should Not Appear One-Sided}{Eliezer Yudkowsky}{2007}

Robin Hanson proposed stores where banned products could be sold. There are good arguments for this: a right to individual liberty, bureaucrats' incentive to ban everything, and legislators as biased as anyone. ``But even so,'' I replied, some poor, honest mother of five would buy ``Dr. Snakeoil's Sulfuric Acid Drink'' for her arthritis and die. \nb{The footnote cites a debate published three months after the post, which does not mention such stores. The proposal was in Hanson's blog post ``Paternalism Is About Bias,'' the day before, and there the buyers had to pass a test showing they understood that regulators disapproved.}

``I was just making a factual observation.'' So why did some people take it as an argument for regulation? My own reply, as I report it, had begun ``But even so,'' which reads as an objection.

On a question of fact, the evidence should end up on one side, because strong evidence is evidence you expect to find on one side only. A policy with many consequences has no reason to be one-sided. People want it to be because ``Arguments are soldiers,'' and the costs of your own policy are enemy soldiers. \nb{These sentences repeat ``Politics is the Mind-Killer'' almost word for word, and again come without evidence.} I also warn against the reverse mistake, splitting the difference between the two positions that get the most airtime. Some policies really are lopsided. But there is ``a human tendency to deny all costs of a favored policy.'' \nb{Ten days later the author cited a study showing that people judge a technology's risks as lower when they judge its benefits as higher. That is a lean, not a denial of all costs.}

The mother is still going to die, whatever you think of the policy. You may also believe that bans only raise prices, that regulators abuse their power, or that her freedom outweighs your wish to meddle; as a matter of fact, she still dies. People dislike an unfair universe, so they deny either the facts or the unfairness. Some libertarians might say that if she walked past signs reading THINGS IN THIS STORE MAY KILL YOU, she deserved it, and then there would be ``no downside'' to the shops. \nb{The post's own first paragraph mentions a second cost: ``her orphans.''} Others believe regulators can be trained to choose well, and then regulation would have no downside.

There is a birth lottery for intelligence. The evidence for a genetic component of 0.6 to 0.8 is ``overwhelming.'' \nb{Summaries of twin and family studies give 0.4 to 0.8, lower in children. The post's point does not depend on the number.} Upbringing is a lottery too. I was raised to think that denying reality is wrong; someone told by a witch doctor that faith is right and doubt is wrong follows the advice in good faith and dies. Would you have been a proper skeptic had you been born in 500 CE?

Real tough-mindedness says the mother did not deserve her death, and keeps the shops open anyway ``because we did this cost-benefit calculation.'' I cannot imagine a politician saying that, but economists who influence policy might at least think it. I count such deaths as tragedies, and draw my moral line at capital punishment, because the dead cannot learn from their mistakes. ``Unfortunately the universe doesn't agree with me.''
